\chapter{Теоретический раздел}


\section{Содержание работы}

Цель работы: изучение венгерского метода решения задачи о назначениях.

\begin{enumerate}
\item реализовать венгерский метод решения задачи о назначениях в виде программы на ЭВМ;
\item провести решение задачи с матрицей стоимостей, заданной в индивидуальном варианте, рассмотрев два случая:
\begin{enumerate}
\item задача о назначениях является задачей минимизации,
\item задача о назначениях является задачей максимизации.
\end{enumerate}
\end{enumerate}

\section{Постановка задачи}

В распоряжении работодателя имеется $n$ работ и $n$ исполнителей. Стоимость выполнения $i$-ой работы $j$-ым исполнителем составляет $c_{ij} \geq 0$ единиц. Требуется распределить работу между исполнителями так, чтобы:
\begin{enumerate}
\item каждый исполнитель выполнял ровно одну работу;
\item общая стоимость выполнения всех работ была минимальной.
\end{enumerate}

\subsection{Математическая формулировка}

Обозначим $C=(c_{ij})_{i,j = \overline{1;n}}$, где $C$ -- матрица стоимостей.
Введем управляемые переменные:

\begin{equation}
x_{ij}(x) =
 \begin{cases}
   1, &\text{если $i$-ую работу выполняет $j$-ый работник}
   \\
   0, &\text{иначе}
 \end{cases},  
\end{equation}
где $i,j = \overline{1,n}$\\

Обозначим $X=(x_{ij})_{i,j = \overline{1;n}}$, где $X$ -- матрица назначений.

Тогда общая стоимость выполнения работ:

\begin{equation}
f=\sum\limits_{i=1}^n\sum\limits_{j=1}^n c_{ij} \cdot x_{ij}
\end{equation}

Условие того, что $j$-й работник выполняет ровно одну работу:

\begin{equation}
\sum\limits_{i=1}^n x_{ij} = 1,
\end{equation}
где $j = \overline{1;n}$

Условие того, что $i$-ю работу выполняет один работник:

\begin{equation}
\sum\limits_{j=1}^n x_{ij} = 1,
\end{equation}
где $i = \overline{1;n}$

Тогда математическая постановка задачи о назначениях имеет следующий вид:

\begin{equation}
x_{ij}(x) =
 \begin{cases}
   f=\sum\limits_{i=1}^n\sum\limits_{j=1}^n c_{ij} \cdot x_{ij} \rightarrow \min
   \\
   \sum\limits_{i=1}^n x_{ij} = 1, \quad j = \overline{1;n}
   \\
   \sum\limits_{j=1}^n x_{ij} = 1, \quad i = \overline{1;n}
   \\
   x_{i,j}\subseteq{0,1}, \quad i,j = \overline{1;n}
 \end{cases},  
\end{equation}
где $f$ -- целевая функция, $c_{ij}$ -- элементы матрицы стоимостей $C$, $x_{ij}$ -- элементы матрицы назначений $X$.


\subsection{Исходные данные индивидуального варианта}

\begin{equation*}
    C = 
\begin{bmatrix}
3 & 5 & 2 & 4 & 8\\
10 & 10 & 4 & 3 & 6\\
5 & 6 & 9 & 8 & 3\\
6 & 2 & 5 & 8 & 4\\
5 & 4 & 8 & 9 & 3
\end{bmatrix}
\end{equation*}


\section{Описание венгерского метода}

Венгерский метод решения задачи о назначениях применяется для решения задачи минимизации, то есть для случая, когда матрица стоимостей $C$ содержит в себе стоимости работ. Также матрица стоимостей $C$ может использоваться для решения задачи максимизации, в этом случае ее элементы отражают прибыль работодателя от назначения $j$-го работника на $i$-ю работу. 

Решение задачи максимизации отличается от решения задачи минимизации тем, что в данном случае в начале алгоритма добавляется шаг, на котором необходимо найти максимальный элемент столбца, вычесть его из каждого элемента матрицы и домножить матрицу на $-1$.

Если решается задача минимизации, алгоритм начинается со следующих шагов:
\begin{enumerate}
\item в каждом столбце матрицы стоимостей выбирается наименьший элемент и вычитается из всех элементов этого столбца;
\item в каждой строке полученной матрицы выбирается наименьший элемент и вычитается из всех элементов этой строки.
\end{enumerate}

В результате получается матрица, в каждой строке и столбце которой присутствует хотя бы один нуль.

Системой независимых нулей (СНН) будем называть такой набор нулей матрицы стоимостей, в котором никакие два элемента не стоят ни в общей строке, ни в общем столбце.

Для построения начальной СНН необходимо просмотреть столбцы текущей матрицы с размерностью $n \times n$ и первый найденный в столбце 0, в одной строке с которым нет элемента $0^*$, отметить как $0^*$. 

Если полученная СНН содержит $n$ элементов, значит, найдено оптимальное решение. Иначе СНН необходимо улучшить.

Для улучшения СНН необходимо:
\begin{enumerate}
\item отметить + столбцы, в которых стоят $0^*$ (данные столбцы и все их элементы будем называть выделенными);
\item найти 0 среди невыделенных элементов и отметить его $0'$;
\item если в одной строке с $0'$ стоит $0^*$, снять выеделение со столбца с найденным $0^*$ и отметить строку с текущим $0'$, затем повторить поиск;

\item если в строке с выделенным $0'$ нет эелемента $0^*$, то строится L-цепочка:

$текущий 0' \rightarrow \text{по столбцу} \rightarrow 0^* \rightarrow \text{по строке} \rightarrow 0' \rightarrow ... \rightarrow 0'$

В пределах цепочки заменить элементы: $0^* \rightarrow 0, 0' \rightarrow 0^*$
\end{enumerate}

Если среди невыделенных элементов не найдено нулей, то необходимо:
\begin{enumerate}
\item среди всех невыделенных элементов выбрать наименьший элемент $h>0$;
\item вычесть h из элементов в невыделенных столбцах;
\item прибавить h к элементам в выделенных строках.
\end{enumerate}

В результате преобразований получим эквивалентную матрицу, среди невыделенных элементов которой есть 0, после этого приступить в улучшению СНН.




% \includeimage
%     {tux} % Имя файла без расширения (файл должен быть расположен в директории inc/img/)
%     {f} % Обтекание (без обтекания)
%     {h} % Положение рисунка (см. figure из пакета float)
%     {0.25\textwidth} % Ширина рисунка
%     {Символ Linux (Tux)} % Подпись рисунка

